\documentclass{article}

\usepackage{corl_2026} % Use this for the initial submission (anonymous, double-blind).
%\usepackage[final]{corl_2026} % Uncomment for the camera-ready ``final'' version.
%\usepackage[preprint]{corl_2026} % Uncomment for pre-prints (e.g., arxiv); This is like ``final'', but will remove the CORL footnote.

% Shared math macros generated by /paper-write.
% Comment out if math_commands.tex does not exist in the project.
\input{math_commands.tex}

\title{Paper Title}

% The \author macro works with any number of authors. There are two
% commands used to separate the names and addresses of multiple
% authors: \And and \AND.
%
% Using \And between authors leaves it to LaTeX to determine where to
% break the lines. Using \AND forces a line break at that point. So,
% if LaTeX puts 3 of 4 authors names on the first line, and the last
% on the second line, try using \AND instead of \And before the third
% author name.

% NOTE: authors will be visible only in the camera-ready and preprint versions (i.e., when using the option 'final' or 'preprint').
% 	For the initial submission the authors will be anonymized.

\author{
  Anonymous Authors\\
  Anonymous Institution\\
  \texttt{anonymous@example.com}
  %% examples of more authors
  %% \And
  %% Coauthor \\
  %% Affiliation \\
  %% Address \\
  %% \texttt{email} \\
  %% \AND
  %% Coauthor \\
  %% Affiliation \\
  %% Address \\
  %% \texttt{email} \\
  %% \And
  %% Coauthor \\
  %% Affiliation \\
  %% Address \\
  %% \texttt{email} \\
  %% \And
  %% Coauthor \\
  %% Affiliation \\
  %% Address \\
  %% \texttt{email} \\
}


\begin{document}
\maketitle

%===============================================================================

\begin{abstract}
    \input{sections/0_abstract.tex}
    % CoRL requires a single paragraph, 4--6 sentences. Gross violations trigger corrections at the camera-ready phase.
\end{abstract}

% Two or three meaningful keywords should be added here
\keywords{Robot Learning, Keyword2, Keyword3}

%===============================================================================

\section{Introduction}
\label{sec:intro}
\input{sections/1_introduction.tex}

%===============================================================================

\section{Related Work}
\label{sec:related}
\input{sections/2_related_work.tex}

%===============================================================================

\section{Method}
\label{sec:method}
\input{sections/3_method.tex}

%===============================================================================

\section{Experiments}
\label{sec:experiments}
\input{sections/4_experiments.tex}

%===============================================================================

\section{Conclusion}
\label{sec:conclusion}
\input{sections/5_conclusion.tex}

%===============================================================================

\clearpage
% The acknowledgments are automatically included only in the final and preprint versions of the paper.
\acknowledgments{If a paper is accepted, the final camera-ready version will (and probably should) include acknowledgments. All acknowledgments go at the end of the paper, including thanks to reviewers who gave useful comments, to colleagues who contributed to the ideas, and to funding agencies and corporate sponsors that provided financial support.}

%===============================================================================

% no \bibliographystyle is required, since the corl style is automatically used.
\bibliography{references}  % .bib

\appendix
\input{sections/A_appendix.tex}

\end{document}
